\section{System Architecture}
\label{sec:implementation}
\begin{figure}[t]
\centering
\begin{tikzpicture}[
  box/.style={draw=black!65, rounded corners=2pt, fill=black!3,
    align=center, font=\footnotesize, minimum height=0.85cm, text width=8.0cm},
  branch/.style={box,text width=6.0cm,minimum height=1.1cm},
  flow/.style={-{Stealth[length=2mm]},thick,draw=black!70},
  note/.style={font=\scriptsize,align=center,fill=white,inner sep=2pt}
]
\node[box] (request) at (0,0)
  {\textbf{Request validation}\\State, typed questions, candidates, optional images};
\node[box] (prepare) at (0,-2.0)
  {\textbf{Context and candidate preparation}\\Semantic rendering, chat template, tokenization\\Optional image encoding and visual projection};
\node[box] (prefix) at (0,-4.0)
  {\textbf{Shared-prefix prefill}\\Retain question state; score candidate first tokens};
\node[branch] (kv) at (-3.5,-6.2)
  {\textbf{KV-cache text models}\\Fork prefix state; batch candidate probes\\Partition into waves when required};
\node[branch] (serial) at (3.5,-6.2)
  {\textbf{Hybrid/recurrent or vision path}\\Preserve prefix; reset working sequence\\Evaluate candidate continuations serially};
\node[box] (score) at (0,-8.5)
  {\textbf{Likelihood aggregation}\\Extract target-token log-probabilities\\Sum per candidate; normalize over candidates};
\node[box] (answer) at (0,-10.5)
  {\textbf{Typed response}\\Choice selection or probability aggregation\\Probabilities, usage, and sequence cleanup};
\draw[flow] (request) -- (prepare);
\draw[flow] (prepare) -- (prefix);
\draw[flow] (prefix.south) -- ++(0,-0.32) -| (kv.north);
\draw[flow] (prefix.south) -- ++(0,-0.32) -| (serial.north);
\draw[flow] (kv.south) -- ++(0,-0.35) -| (score.north);
\draw[flow] (serial.south) -- ++(0,-0.35) -| (score.north);
\draw[flow] (score) -- (answer);
\end{tikzpicture}
\caption{End-to-end sampler-free workflow, executed locally on CPU, GPU, or a
mixed placement. Candidate tokens are supplied rather than sampled. Optional
text-only reasoning, omitted from the principal path above, generates a trace
before scoring and freezes it into the shared prefix; that mode includes
sampling. Invalid requests produce errors.}
\label{fig:workflow}
\end{figure}

\subsection{End-to-End Execution}
The workspace separates protocol and normalization utilities (\texttt{jet-core}),
engine orchestration (\texttt{jet-engine}), native bindings
(\texttt{jet-llama-sys}), and the command-line interface (\texttt{jet-cli}).
Figure~\ref{fig:workflow} summarizes the local execution workflow. Each question
passes through semantic rendering, template application,
tokenization, prefix execution, candidate evaluation, and typed answer assembly.
The Vulkan and CUDA backends are developed on top of llama.cpp and its ggml
execution backends~\cite{llamacpp}. JET builds and links the corresponding
backends from a pinned llama.cpp submodule. Model loading, quantized tensor
operations, and native GPU kernels are provided by this upstream runtime; JET
adds decision-specific scoring, prefix-state management, candidate scheduling,
and host-side orchestration. The default build supports CPU execution; CUDA,
Metal, Vulkan, and vision are optional build features. The measurements below concern CPU and Vulkan paths,
not a comparative study of all supported backends.

\subsection{Shared-Prefix Candidate Probing}
We use \emph{probe} to mean evaluating a supplied candidate continuation, not
training a probing classifier. Let a question have $K$ candidates and a prefix
of $P$ tokens. Independent evaluation processes that prefix $K$ times. If all
candidates can reuse one prefix state, the number of prefix tokens processed
falls from $KP$ to $P$. This token-count argument does not imply a proportional
wall-time reduction: model computation, memory movement, candidate lengths, and
batch geometry also contribute.

For pure KV-cache models, JET prefills a question once when its candidates fit
the sequence budget, copies the sequence memory, and evaluates candidate
continuations in native batches. Larger workloads run in waves without
truncating the question. All candidate tokens are known, permitting
teacher-forced continuation evaluation rather than repeated sampled generation.
The phrase parallel probing describes this supported batch schedule; it does
not imply unlimited simultaneous candidates or a single forward pass for every
request size.

Hybrid recurrent models require a different schedule. Sequence 0 holds the
unchanged prefix. Its final output scores each candidate's first token. Before
advancing a remaining continuation, JET resets sequence 1 and copies the prefix
state into it. Only that candidate sequence advances; shared recurrent state is
preserved through the native backend's copy-on-write behavior. Two sequence slots
suffice, including when there are more than two candidates. A one-token answer
requires no continuation decode. This retains prefix reuse while avoiding the
incorrect sibling-state interaction observed in an earlier batched schedule.

After a scoring wave or serial question completes, JET releases sequence
metadata instead of synchronously zeroing the entire allocated cache. Attention
masks exclude old KV entries and the backend initializes fresh recurrent state
when decoding a new sequence. Initial native context allocation still performs
its normal buffer initialization.

\subsection{Probability Extraction and Preparation Overlap}
The batched GPU scoring path performs full-vocabulary softmax, target-token
gather, and log computation on the backend holding the output weights. It returns
target values rather than full vocabulary rows. The gather width can grow between
questions to cover distinct first tokens independently of the sequence-slot
count. This reduces transfer and buffer requirements without changing the
candidate definition.

Native chat-template objects and tokenizer buffers are retained across requests.
For disabled-thinking, batched Vulkan hybrid execution, a persistent CPU worker
prepares upcoming questions behind a bounded queue while the GPU scores the
current question. Results and errors remain in input order. Preparation overlap
does not make recurrent candidate execution parallel, and
\texttt{--batch-requests} controls the input chunk rather than the native
inference batch. Synchronous execution remains available for comparison.

\subsection{Placement on Consumer Hardware}
GPU layer placement and CPU expert placement are explicit configuration choices.
For models larger than available VRAM, routed experts can remain on CPU while
attention, shared experts, and other layers execute on GPU. These are computation
placements, not per-token streaming of selected expert weights from disk.
A scorer-owned sleeping CPU worker pool persists across graph segments and is
shared by the sequential scoring and thinking contexts. The mixed CPU/Vulkan
results in Section~\ref{sec:offload} show why worker lifetime matters when CPU
expert segments repeatedly interrupt GPU execution.

\subsection{Multimodal Inference}
\label{sec:multimodal}
The same scoring rule can condition textual candidates on both images and text
using a vision-language model and its matching visual projector. Visual
embeddings count toward the context budget. The current implementation requires
reasoning to be disabled.

A Doom controller illustrates this extension. It selects among nine admissible
controls using 160-by-100 RGB screenshots, visible status information, and recent
actions. Hidden enemy positions and map geometry are excluded. Visual decision
accuracy and episode-level control performance have not been quantitatively
evaluated.

